\documentclass[11pt,a4paper]{article}

\usepackage[utf8]{inputenc}
\usepackage[T1]{fontenc}
\usepackage[spanish]{babel}
\usepackage{amsmath,amssymb,mathtools}
\usepackage{geometry}
\usepackage{enumitem}
\usepackage{array}
\usepackage{lmodern}
\usepackage{microtype}

\geometry{margin=2.4cm}
\setlength{\parindent}{0pt}
\setlength{\parskip}{0.5em}

\newcommand{\R}{\mathbb{R}}
\newcommand{\K}{\mathbb{K}}
\newcommand{\Ker}{\operatorname{Ker}}
\newcommand{\Ima}{\operatorname{Im}}
\newcommand{\rg}{\operatorname{rg}}
\newcommand{\Span}{\operatorname{L}}
\newcommand{\Mat}{\operatorname{M}}
\newcommand{\Id}{\operatorname{Id}}
\newcommand{\an}{\operatorname{an}}

\begin{document}

\begin{center}
	{\LARGE\bfseries Álgebra Lineal}\\[0.4em]
	{\Large Tema 2: Aplicaciones Lineales --- Ejercicios}
\end{center}

\section{Ejercicio 1}
Estudiar si es lineal:

\[
f:\R^3\longrightarrow\R^4
\]

\[
f(x_1,x_2,x_3)
=
(x_3,x_1+x_2,0,-x_3)
\]

\section{Ejercicio 2}
Estudiar si es lineal:

\[
f:\R^3\longrightarrow\R^2
\]

\[
f(x,y,z)=(x+a,y)
\]

\section{Ejercicio 3}
Estudiar si es lineal:

\[
f:\R^3\longrightarrow M_{2\times2}(\R)
\]

\[
f(x,y,z)
=
\begin{pmatrix}
ax&by\\
cz&x+y+z
\end{pmatrix},
\qquad
a,b,c\in\R
\]

\section{Ejercicio 4}
Dado el endomorfismo de $\R^3$:

\[
f(x_1,x_2,x_3)
=
(x_1+x_2+x_3,\;x_1+x_2-x_3,\;x_3)
\]

obtener:

\begin{enumerate}[label=\alph*)]
	\item $\Ker f$.
	\item $\Ima f$.
	\item $f(V)$, siendo:
	\[
	V=
	\{(x_1,x_2,x_3)\in\R^3:x_1+x_2+x_3=0\}
	\]
\end{enumerate}

\section{Ejercicio 5}
Dada:

\[
f:\R^3\longrightarrow\R^2
\]

definida por:

\[
f(x_1,x_2,x_3)
=
(x_1+2x_2+3x_3,\;2x_1+4x_2+6x_3)
\]

obtener:

\begin{enumerate}[label=\alph*)]
	\item $\Ker f$.
	\item $\Ima f$.
	\item Comprobar sus dimensiones.
\end{enumerate}

\section{Ejercicio 6}
Determinar $\Ker f$ e $\Ima f$ del endomorfismo:

\[
f:\R^3\longrightarrow\R^3
\]

\[
f(x_1,x_2,x_3)
=
(x_3,x_1+x_2,-x_3)
\]

y calcular:

\[
f^{-1}(0,0,0),
\qquad
f^{-1}(1,2,3),
\qquad
f^{-1}(1,0,-1)
\]

\section{Ejercicio 7}
Dada la aplicación lineal:

\[
f:\mathcal{P}_2(x)\longrightarrow\mathcal{P}_2(x)
\]

tal que:

\[
f(1+x^2)=3+2x^2
\]

\[
f(x)=1+x+2x^2
\]

\[
f(3+x)=-1
\]

determinar la imagen del polinomio:

\[
1-x+5x^2
\]

\section{Ejercicio 8}
Sean:

\[
B=\{u_1,u_2\}
\]

y:

\[
B'=\{v_1,v_2,v_3\}
\]

bases de $U$ y $V$, respectivamente. Si $g:U\to V$ es lineal y:

\[
g(3u_1-2u_2)=3v_1+6v_2-3v_3
\]

\[
g(4u_1-3u_2)=v_1+5v_2-v_3
\]

obtener:

\[
M_{B,B'}(g)
\]

y clasificar el homomorfismo.

\section{Ejercicio 9}
Sea:

\[
f:\R^3\longrightarrow\R^2
\]

\[
f(x_1,x_2,x_3)
=
(x_1+3x_2+5x_3,\;2x_1+2x_2+7x_3)
\]

Calcular la matriz asociada a $f$ en las bases:

\[
A'=
\{(1,1,1),(1,1,0),(0,1,1)\}
\]

\[
B'=
\{(1,1),(4,3)\}
\]

\section{Ejercicio 10}
Sean:

\[
A'=\{u_1,u_2,u_3\}
\]

una base de $U$ y:

\[
B'=\{v_1,v_2\}
\]

una base de $V$. Sea $f:U\to V$ definida por:

\[
f(u_1+3u_2+2u_3)=v_1+3v_2
\]

\[
f(u_2+u_3)=-v_1+v_2
\]

\[
f(u_1+u_3)=4v_1+2v_2
\]

Calcular:

\[
M_{A',B'}(f)
\]

\section{Ejercicio 11}
Sea $C$ la base canónica de $\R^3$ y:

\[
C'=
\{(1,1,0),(1,0,1),(0,-1,2)\}
\]

Si $f$ es un endomorfismo cuya matriz respecto de la base canónica es:

\[
M_C(f)=
\begin{pmatrix}
1&0&0\\
1&0&-1\\
-2&2&3
\end{pmatrix}
\]

calcular la matriz de $f$ en la base $C'$.

\section{Ejercicio 12}
Sea:

\[
f:\R^3\longrightarrow\R^3
\]

lineal. Consideramos:

\[
S=
\{(x,y,z)\in\R^3:x+y+z=0\}
\]

y se sabe que:

\[
f(v)=v
\qquad
\forall v\in S
\]

y:

\[
\Ker f=
\{w\in\R^3:w^tv=0\quad\forall v\in S\}
\]

Se pide:

\begin{enumerate}[label=\alph*)]
	\item Determinar $M_{B,B_c}(f)$, donde:
	\[
	B=\{(1,0,-1),(0,1,-1),(1,1,1)\}
	\]
	\item Dar la expresión analítica de $f(x,y,z)$ en bases canónicas.
	\item Calcular $f^{-1}(W)$, siendo:
	\[
	W=\{(x,y,z)\in\R^3:y=z=0\}
	\]
	\item Razonar las dimensiones de $\Ker f$ e $\Ima f$ y clasificar $f$.
\end{enumerate}

\section{Ejercicio 13}
Sea:

\[
f:\R^3\longrightarrow\R^2
\]

con:

\[
M_{B_c,B_c}(f)
=
\begin{pmatrix}
3&0&-2\\
-1&4&5
\end{pmatrix}
\]

Calcular la matriz asociada a $f$ en las bases:

\[
B=
\{(1,3,0),(1,0,2),(0,4,-2)\}
\]

\[
B'=
\{(2,1),(4,3)\}
\]

\section{Ejercicio 14}
Sea $S_2$ el espacio de matrices simétricas de orden $2$ y:

\[
B_1=
\left\{
\begin{pmatrix}1&0\\0&0\end{pmatrix},
\begin{pmatrix}0&1\\1&0\end{pmatrix},
\begin{pmatrix}0&0\\0&1\end{pmatrix}
\right\}
\]

su base canónica. Sea $B_2$ la base canónica de $\R^3$.

Sean:

\[
f:S_2\longrightarrow\R^3,
\qquad
g:\R^3\longrightarrow S_2
\]

aplicaciones lineales tales que:

\[
\Ker f=
\{(x,y,z)_{B_1}:x-y+2z=0\}
\]

y:

\[
(g\circ f)
\begin{pmatrix}
1&1\\
1&1
\end{pmatrix}
=
\begin{pmatrix}
-1&0\\
0&2
\end{pmatrix}
\]

Sea:

\[
h=g\circ f
\]

Se pide:

\begin{enumerate}[label=\alph*)]
	\item Calcular la matriz de $h$ en la base $B_1$.
	\item Obtener una base de $h(W)$, donde:
	\[
	W=
	\left\{
	\begin{pmatrix}
	x&y\\
	y&z
	\end{pmatrix}:
	\begin{array}{rcl}
	x-2y+3z&=&0\\
	3x-7y+7z&=&0\\
	5x-11y+13z&=&0
	\end{array}
	\right\}
	\]
	\item Si:
	\[
	B'=
	\left\{
	\begin{pmatrix}1&0\\0&2\end{pmatrix},
	\begin{pmatrix}0&1\\1&-1\end{pmatrix},
	\begin{pmatrix}-3&1\\1&1\end{pmatrix}
	\right\}
	\]
	es otra base de $S_2$, obtener la matriz de $h$ en $B'$.
	\item Razonar si:
	\[
	x+2y+3z=0,
	\qquad
	y-4z=0
	\]
	pueden ser ecuaciones implícitas en $B_1$ del subespacio $\Ima g$.
\end{enumerate}

\section{Ejercicio 15}
Dada la base:

\[
B=
\{(1,1,1),(1,1,0),(1,0,0)\}
\]

de $\R^3$, obtener la base dual.

\section{Ejercicio 16}
Sea:

\[
B=
\{v_1=(1,2),v_2=(3,1)\}
\]

una base de $\R^2$ y:

\[
B^*=\{v_1^*,v_2^*\}
\]

su base dual. Sea la forma lineal:

\[
f=3v_1^*+4v_2^*
\]

Determinar $f(x_1,x_2)$ para un vector cualquiera de $\R^2$.

\section{Ejercicio 17}
Sea:

\[
f_\beta:\mathcal{P}_2(x)\longrightarrow M_{2\times2}(\R)
\]

definida por:

\[
f_\beta(p)
=
\begin{pmatrix}
\beta p(1)&p'(1)\\
p''(1)&p'''(1)
\end{pmatrix}
\]

Se consideran las bases:

\[
B_1=\{1,x,x^2\}
\]

y la base canónica $B_2$ de $M_{2\times2}(\R)$.

Se pide:

\begin{enumerate}[label=\alph*)]
	\item Obtener $M_{B_1,B_2}(f_\beta)$.
	\item Calcular $\dim(\Ima f_\beta)$ y $\dim(\Ker f_\beta)$ según $\beta$.
	\item Para el caso no inyectivo, obtener ecuaciones paramétricas de $\Ker f_\beta$ en $B_1$ y ecuaciones implícitas de $\Ima f_\beta$ en $B_2$.
	\item Para $\beta=1$, hallar una base de:
	\[
	f_1^{-1}(S_{2\times2})
	\]
	siendo $S_{2\times2}$ el espacio de matrices simétricas.
	\item Para $\beta=1$, considerando:
	\[
	B_1'=\{1,x-1,1+x^2\}
	\]
	y:
	\[
	B_2'=
	\left\{
	\begin{pmatrix}1&0\\0&0\end{pmatrix},
	\begin{pmatrix}0&-1\\0&0\end{pmatrix},
	\begin{pmatrix}0&0\\-1&0\end{pmatrix},
	\begin{pmatrix}1&1\\-1&2\end{pmatrix}
	\right\}
	\]
	obtener $M_{B_1',B_2'}(f_1)$.
\end{enumerate}

\section{Ejercicio 18}
En $\mathcal{P}_2(x)$ se define:

\[
f:\mathcal{P}_2(x)\longrightarrow\mathcal{P}_2(x)
\]

tal que:

\begin{enumerate}[label=\arabic*)]
	\item Los polinomios con término independiente nulo se transforman en sí mismos.
	\item
	\[
	\Ker f=
	\{a+bx+cx^2:a=b=c\}
	\]
\end{enumerate}

Se pide:

\begin{enumerate}[label=\alph*)]
	\item Calcular la matriz asociada a $f$ en la base canónica.
	\item Si:
	\[
	S=
	\{
	\alpha-2\beta+\gamma
	+
	(\alpha-2\beta-\gamma)x
	+
	(\alpha-2\beta)x^2
	\}
	\]
	obtener una base de $f(S)$.
	\item Clasificar el homomorfismo.
	\item Utilizando la matriz de $f$ en:
	\[
	B=\{1+x+x^2,1+x,1\}
	\]
	determinar las coordenadas canónicas de:
	\[
	f(1-3x+5x^2)
	\]
\end{enumerate}

\section{Ejercicio 19}
Sea:

\[
B=
\{(1,1,0),(0,1,1),(1,0,1)\}
\]

una base de $\R^3$. Encontrar su base dual.

\section{Ejercicio 20}
Sea $f$ un endomorfismo de $\R^4$ definido por:

\[
\Ker f=
\left\{
(x,y,z,t):
\begin{array}{rcl}
2x+y-z-2t&=&0\\
z+2t&=&0
\end{array}
\right\}
\]

y:

\[
f(0,0,0,1)=(2,0,0,0)
\]

\[
f(1,0,0,0)=(2,0,2,0)
\]

Se pide:

\begin{enumerate}[label=\alph*)]
	\item Calcular la matriz de $f$ respecto de la base canónica.
	\item Hallar una base de $f(V)$, donde:
	\[
	V=
	\{(x,y,z,t):x+y+z+t=0\}
	\]
	\item Calcular la matriz de $f$ respecto de:
	\[
	B=
	\{(1,1,0,0),(1,-1,0,0),(0,0,1,1),(0,0,1,-1)\}
	\]
\end{enumerate}

\section{Ejercicio 21}
Sean $A,B,C$ espacios vectoriales con bases:

\[
B_A=\{e_1,e_2,e_3\},
\qquad
B_B=\{u_1,u_2\},
\qquad
B_C=\{c_1,c_2,c_3\}
\]

y:

\[
f:A\to B
\]

definida por:

\[
f(e_1)=u_1-u_2,
\qquad
f(e_2)=u_2,
\qquad
f(e_3)=2u_2
\]

y:

\[
g:B\to C
\]

definida por:

\[
g(u_1)=c_1-c_2+2c_3,
\qquad
g(u_2)=c_1-c_2
\]

Sea:

\[
h=g\circ f
\]

Calcular:

\begin{enumerate}[label=\alph*)]
	\item La matriz de $h$.
	\item $h^{-1}(1,1,1)$.
	\item $\Ker h$.
	\item $h(V\cap W)$, siendo:
	\[
	V=
	\{(2\alpha+\beta,\alpha-\beta,-\alpha):\alpha,\beta\in\R\}
	\]
	y:
	\[
	W=
	\{(x,y,z):x-y+2z=0\}
	\]
\end{enumerate}

\section{Ejercicio 22}
Sean:

\[
E=\{ax+b:a,b\in\R\}
\]

\[
F=S_2(\R)
\]

\[
G=\R^3
\]

y los homomorfismos:

\[
f:E\to F,
\qquad
f(ax+b)
=
\begin{pmatrix}
a&b\\
b&a
\end{pmatrix}
\]

\[
g:F\to G,
\qquad
g
\begin{pmatrix}
a&d\\
d&c
\end{pmatrix}
=
(a,c,a+c)
\]

Con las bases canónicas:

\[
B=\{x,1\},
\qquad
B',
\qquad
B''
\]

de $E$, $F$ y $G$, respectivamente, calcular:

\begin{enumerate}[label=\alph*)]
	\item Las matrices de $f$, $g$ y $g\circ f$.
	\item $(g\circ f)(V)$, siendo:
	\[
	V=\{ax+a:a\in\R\}
	\]
	\item $\Ker g$ e $\Ima g$ y clasificar $g$.
	\item Si en $F$ se considera:
	\[
	B^*=
	\left\{
	\begin{pmatrix}1&0\\0&1\end{pmatrix},
	\begin{pmatrix}0&1\\1&0\end{pmatrix},
	\begin{pmatrix}1&0\\0&-1\end{pmatrix}
	\right\}
	\]
	hallar de nuevo las matrices de $f$, $g$ y $g\circ f$.
\end{enumerate}

\section{Ejercicio 23}
Sean $S_2(\R)$ y $\mathcal{P}_2(x)$ y se define:

\[
f:\mathcal{P}_2(x)\longrightarrow S_2(\R)
\]

mediante:

\[
f[p(x)]
=
\begin{pmatrix}
p'(0)&p'(1)\\
p'(1)&p'(-1)
\end{pmatrix}
\]

Se pide:

\begin{enumerate}[label=\alph*)]
	\item Probar que $f$ es lineal y obtener su matriz en las bases canónicas.
	\item Calcular las ecuaciones implícitas y una base de $\Ker f$ en:
	\[
	B=\{x^2,(x-1)^2,(x+1)^2\}
	\]
\end{enumerate}

Sea además:

\[
g:S_2(\R)\longrightarrow\mathcal{P}_1(x)
\]

definida por:

\[
g
\begin{pmatrix}
a&b\\
b&d
\end{pmatrix}
=
(a+d)+(a-d)x
\]

y sea:

\[
h=g\circ f
\]

Calcular:

\begin{enumerate}[label=\alph*),start=3]
	\item Una base de $h^{-1}(T)$, siendo:
	\[
	T=\{k+kx:k\neq0\}
	\]
	\item La matriz $M_{B,B'}(h)$, siendo:
	\[
	B'=\{(1,-2),(0,1)\}
	\]
	\item A partir de dicha matriz, calcular:
	\[
	h(3-2x+6x^2)
	\]
	y dar sus coordenadas en la base canónica.
\end{enumerate}

\end{document}
